% Junior Level 1, Week 10 -- answer key
\answersweek{10}{Showcase and Certificate}{Answers for workbook chapter 10 and play-at-home card 10.}

\section{Lesson questions}
\begin{answer}[Exit question]
The action named in the hat block: click the green flag (\enquote{when green flag clicked}), click that sprite (\enquote{when this sprite clicked}), or press the named key (\enquote{when space key pressed}).
\end{answer}

\begin{answer}[Reflection]
Any drawing of a moment from the card and any honest smiley is right. Answers to \enquote{What would you like to make next?} help you plan Junior Level 2 (games), so write them down.
\end{answer}

\begin{answer}[Show-me checklist]
The Week 10 show-me checklist is in the Assessments booklet. It is observed while the child presents and is never graded. \enquote{Getting there} is a normal result at this age.
\end{answer}

\section{Workbook exercises}
\begin{answer}[Exercise 10.1: Last-Minute Bug Check]
A finished card gets \enquote{yes} in all four rows. The most common \enquote{not yet} is row 2: the card starts differently the second time, because a sprite stayed hidden or the backdrop stayed changed. Fix: \enquote{show} or \enquote{switch backdrop to} (purple Looks) at the top of a green flag script. No sound: check the Mac volume first, then the sound block. Fixing one thing is enough today.
\end{answer}

\begin{answer}[Exercise 10.2: My Showcase]
Any true sentence about the child's own card is right, for example \enquote{My card is for Grandpa. When you click the balloon, it pops and says Surprise!} or \enquote{When you press space, the backdrop changes.} Look for: the child names what starts the script (the hat block) and one thing that happens next. Pointing at the right script while saying it counts as \enquote{Got it} for explaining.
\end{answer}

\begin{answer}[Exercise 10.3: Sticker Feedback]
Any sticker placed with a kind spoken reason is right (\enquote{I liked the sound, because it made me laugh}). Look for: the visitor asks before clicking, uses kind words only, and does not change the card.
\end{answer}

\section{Play-at-home card}
\begin{answer}[Show your card at home]
Look for the same showcase sentence as in class: what starts a script and what it does. The family member's favorite part can be anything.
\end{answer}

\begin{answer}[What next?]
Any idea is right. Games, stories and pen drawings all lead into Junior Level 2 (Game Makers).
\end{answer}
